\documentclass[11pt]{article}
\usepackage[margin=1in]{geometry}
\usepackage{amsmath,amssymb}
\title{The Galois group of the \(c_1\)-spectrum of \(\mathrm{Gr}(3,8)\)}
\author{The Last Math Competition, conjecture 00000002104}
\date{}
\begin{document}
\maketitle

Let \(P_q\) be the characteristic polynomial of multiplication by
\(c_1(T\mathrm{Gr}(3,8))\) in small quantum cohomology.
We compute its Galois group both over \(\mathbb Q\) at \(q=1\)
and over \(\mathbb Q(q)\) with \(q\) indeterminate.

\paragraph{Spectrum.}
Let \(S\) be the rank-three tautological bundle and \(Q\) the
rank-five quotient bundle. The tangent bundle is \(S^*\otimes Q\),
and \(c_1(Q)=c_1(S^*)=\sigma_1\). Therefore \(c_1(T\mathrm{Gr}(3,8))=8\sigma_1\).
Under the standard presentation of quantum cohomology, \(\sigma_1\)
is the first elementary symmetric function \(e_1\).
The geometric points of the nonzero \(q\)-fiber are indexed by the
three-element subsets of the roots of \(z^8=q\), and \(e_1\) evaluates
there as the sum of the three elements
\cite[Lemma 3.7 and Theorem 3.4]{rietsch}.
For completeness, the presentation over \(\mathbb Q[q]\) is
\[
 \mathbb Q[q,e_1,e_2,e_3]/(h_6,h_7,h_8-q),
 \qquad
 \sum_{m\geq0}h_m t^m=(1-e_1t+e_2t^2-e_3t^3)^{-1},
\]
by the Siebert--Tian theorem \cite[Theorem 1]{galkin}.
For a three-element subset of roots of \(z^8=q\), the generating
function is the product of \(1-zt\) over the five complementary
roots divided by \(1-qt^8\). Thus these evaluations satisfy the
relations. They give distinct maximal ideals after extending the
ground field to contain all the roots of \(z^8=q\), since the three
elementary symmetric coordinates determine the subset. The
Chinese remainder theorem gives a surjection onto the product of
the resulting \(\binom83=56\) fields, and the quantum ring has
dimension \(56\) by its Schubert basis. This proves that the
evaluations account for the entire characteristic polynomial.

Fix \(\alpha^8=q\) and a primitive eighth root of unity \(\zeta\).
Writing \(r_{ijk}\) for the eigenvalue indexed by
\(\{i,j,k\}\subset\{0,\ldots,7\}\), we have
\[
 r_{ijk}=8\alpha(\zeta^i+\zeta^j+\zeta^k).
\]
The 56 subset labels give the eigenvalues with multiplicity.

\paragraph{At \(q=1\).}
Choose \(\alpha=1\). All eigenvalues lie in \(\mathbb Q(\zeta)\), and
\[
 \zeta=1+\frac{r_{123}-r_{023}}{8}
\]
lies in their splitting field. Thus the splitting field is exactly
\(\mathbb Q(\zeta)\). The polynomial \(x^4+1\) is irreducible over
\(\mathbb Q\), since replacing \(x\) by \(x+1\) gives a polynomial
Eisenstein at \(2\). Its roots are the odd powers of \(\zeta\).
The Galois group is therefore
\((\mathbb Z/8\mathbb Z)^\times\cong C_2\times C_2\), of order \(4\).

\paragraph{Over \(\mathbb Q(q)\).}
Every eigenvalue lies in \(\mathbb Q(q,\zeta,\alpha)\). Since
\(\zeta^4=-1\), the reverse inclusion follows from
\[
 \alpha=\frac{r_{023}-r_{234}}{16},
 \qquad
 \zeta=1+\frac{r_{123}-r_{023}}{8\alpha}.
\]
Here \(q\ne0\), so \(\alpha\ne0\). The cyclotomic constant field
extension \(\mathbb Q(\zeta)(q)/\mathbb Q(q)\) has degree \(4\),
because adjoining an indeterminate preserves the degree of a finite
constant extension. The polynomial \(x^8-q\) is Eisenstein at \(q\)
in \(\mathbb Q(\zeta)[q][x]\), so the further extension adjoining
\(\alpha\) has degree \(8\). The characteristic polynomial is
separable after removing repeated factors, so its splitting field
is Galois and its group has order \(32\).
More explicitly, the 32 substitutions
\(\alpha\mapsto\zeta^r\alpha\), \(\zeta\mapsto\zeta^u\) are determined
by \(r\bmod8\) and \(u\in(\mathbb Z/8\mathbb Z)^\times\).
All are automorphisms since there are 32 possibilities and the
Galois group has order 32. Composition sends
\((r,u),(s,v)\) to \((r+us,uv)\), with entries taken modulo eight.

The group at \(q=1\) has order \(4\) and the generic group has order
\(32\), whereas \(\lvert S_5\rvert=120\). This refutes the asserted
Galois group in both interpretations.

\paragraph{Lean verification.}
The accompanying Lean project models \(\zeta^i\) in the basis
\(1,\zeta,\zeta^2,\zeta^3\), verifies the subset indexing and the two
splitting-field extraction identities, and checks the finite Galois
action, composition law, and group orders. The quantum cohomology
presentation and the field-theoretic identifications are justified
above.

\begin{thebibliography}{2}
\bibitem{rietsch}
K. Rietsch, \emph{Quantum cohomology rings of Grassmannians and total
positivity}, Duke Math. J. \textbf{110} (2001), 523--553.
arXiv:math/0112022.
\bibitem{galkin}
S. Galkin and V. Golyshev, \emph{Quantum cohomology of Grassmannians
and cyclotomic fields}, Russian Math. Surveys \textbf{61} (2006),
171--173. DOI: \texttt{10.1070/RM2006v061n01ABEH004304}.
\end{thebibliography}
\end{document}
